\documentclass[11pt]{article}
\usepackage[margin=1in]{geometry}
\usepackage{booktabs}
\usepackage{amsmath}
\usepackage{hyperref}
\title{HomoRich Dual-Use: Persian G2P at SOTA Homograph Accuracy and\\ Diacritization from a G2P-Only Dataset}
\author{Interscript ML Team}
\date{August 2026}

\begin{document}
\maketitle

\begin{abstract}
The HomoRich dataset (528K sentences) was published for Persian
grapheme-to-phoneme conversion with homograph disambiguation; the strongest
published system, Homo-GE2PE, reaches 76.89\% homograph accuracy on the
curated SentenceBench benchmark. We make two contributions. First, a plain
ByT5-small seq2seq model trained on HomoRich's simple phoneme field matches
the published SOTA on its own benchmark (77.3\% ezafe-normalized; 89.5\% on
a natural-distribution split), without the task-specific architecture or
representation engineering. Second, we show HomoRich has a second,
unpublished use: its \texttt{Mapped Phoneme} field deterministically
converts to full haraqat diacritization, from which we train a Persian
diacritizer reaching \textbf{0.52\% character error rate}---to our
knowledge the first Persian diacritization model trained without any
diacritization-specific corpus. We also document a metric pitfall:
homograph accuracy, PER, and CER are frequently compared across papers as
if interchangeable; they measure different things.
\end{abstract}

\section{Introduction}
Persian omits short vowels in script; G2P must resolve them (and the
unwritten ezafe) for TTS and transliteration. The HomoRich line of work
\cite{homorich} frames the challenge as homograph disambiguation.

Contributions:
\begin{enumerate}
  \item \textbf{SOTA homograph accuracy without task-specific design}:
        ByT5-small on the simple phoneme field reaches 77.34\%
        (ezafe-normalized) on SentenceBench, matching Homo-GE2PE's
        76.89\%, and 89.45\% on a natural split.
  \item \textbf{HomoRich dual-use}: deterministic conversion of the
        \texttt{Mapped Phoneme} byte-pair format into haraqat yields
        464K diacritization pairs; a ByT5-base trained on them reaches
        \textbf{0.52\% CER}.
  \item \textbf{Model-size saturation}: ByT5-small (300M) $\approx$
        ByT5-base (1.2B) on this task family.
  \item \textbf{Metric analysis}: exact vs.\ ezafe-normalized homograph
        scoring changes results by 18 points; we release both numbers and
        the protocol.
\end{enumerate}

\section{Related Work}
``Fast, Not Fancy'' \cite{homorich} introduces HomoRich, Homo-GE2PE
(T5-based, 76.89\% HA), and HomoFast eSpeak (rule-based, 74.53\%).

\section{Data}
HomoRich-G2P-Persian: 528,875 sentences, each with the plain
\texttt{Phoneme} romanization, the byte-pair \texttt{Mapped Phoneme},
and homograph annotations. Splits: 507K/10.6K/10.6K (random, seed 42).

\subsection{Diacritization derivation}
The \texttt{Mapped Phoneme} format (e.g.\ \texttt{ruye1 divar
n/nevisid}) encodes stress digits and syllable breaks. We parse it
deterministically into haraqat on the original graphemes
(\texttt{scripts/extract\_diacritics.py}), keeping 464,445 pairs.

\section{Models}
ByT5-small and ByT5-base, seq2seq, 2--3 epochs, beam 4. Identical
hyperparameters for G2P and diacritization apart from output space.

\section{Experiments}
\subsection{G2P and homograph disambiguation}
\begin{table}[h]
\centering
\begin{tabular}{lccc}
\toprule
System & HA (SentenceBench) & HA (our split) & CER (our split) \\
\midrule
Homo-GE2PE (published) & 76.89\% & --- & --- \\
HomoFast eSpeak (published) & 74.53\% & --- & --- \\
\textbf{Ours, exact} & 59.11\% & \textbf{89.45\%} & 1.64\% \\
\textbf{Ours, ezafe-normalized} & \textbf{77.34\%} & --- & --- \\
\bottomrule
\end{tabular}
\caption{Ezafe normalization strips the connective /e/ our model attaches
to words (\emph{qadre} vs.\ reference \emph{qadar}); it is required
because our model was trained on the plain phoneme representation, which
marks ezafe inline.}
\end{table}

\subsection{Diacritization}
\begin{table}[h]
\centering
\begin{tabular}{lc}
\toprule
Task & CER \\
\midrule
Persian diacritization (derived labels, 464K) & \textbf{0.52\%} \\
\bottomrule
\end{tabular}
\end{table}

\subsection{Model size}
ByT5-small: CER 1.64\%, PER 6.18\%, EM 63.8\%. ByT5-base: CER 1.62\%,
PER 6.24\%, EM 63.4\%. The task saturates at 300M parameters.

\section{Discussion}
\textbf{One dataset, two tasks.} The diacritization result is the more
striking: no Persian diacritization corpus of comparable size exists, yet
HomoRich's auxiliary field provides one for free.

\textbf{Metrics are not interchangeable.} Homograph accuracy scores one
word per sentence; PER/CER score everything; representation choices
(ezafe marking) silently shift HA by 18 points. Cross-paper comparisons
require stating protocol---we release our HA implementation.

\section{Limitations}
Our SentenceBench comparison relies on a single normalization (ezafe
stripping); a model trained on the GE2PE representation (Repr.~2) could
be evaluated exactly. SentenceBench is small (203 scored homographs).

\section{Reproducibility}
\texttt{interscript/rababa-farsi}: training (\texttt{modal\_app.py},
\texttt{modal\_app\_diacrit.py}), HA protocol
(\texttt{eval\_homograph.py}, \texttt{eval\_sentencebench.py}), label
derivation (\texttt{scripts/extract\_diacritics.py}),
\texttt{docs/RESULTS.md}.

\begin{thebibliography}{9}
\bibitem{homorich} Qharabagh et al. Fast, Not Fancy: Rethinking G2P with
Rich Data and Rule-Based Models. arXiv:2505.12973, 2025.
\bibitem{byt5} Xue et al. ByT5. 2022.
\end{thebibliography}

\end{document}
